\section{Introduction}\label{sec:ntro}

In this paper we present a simple mathematical theory for modeling
streaming and incremental computations.  This model has immediate
practical applications in the design and implementation of streaming databases
and incremental view maintenance.  Our model is based on mathematical formalisms 
used in discrete digital signal processing (DSP)~\cite{rabiner-book75},
but we apply it to database computations.  
Thus, we have called it ``\dbsp''.  \dbsp is inspired from Differential 
Dataflow~\cite{mcsherry-cidr13} (DD), and started as an attempt to provide a simpler
formalization of DD than the one of Abadi et al.~\cite{abadi-fossacs15} 
(as discussed in \secref{sec:related}), but has evolved behind that purpose.

The core concept of \dbsp is the \emph{stream}: a stream $s$ with type 
$\stream{A}$ maps ``time'' moments $t\in\N$ 
to values $s[t]$ of type $A$; think of it as an "infinite vector".
A streaming computation is a function that
consumes one or more streams and produces another stream.  We depict
streaming computations with typical DSP box-and-arrow diagrams (also called ``circuits''),
where boxes are computations and streams are arrows, as in the following diagram,
which shows a stream operator $T$ consuming two input streams $s_0$ and $s_1$ 
and producing one output stream $s$:

\begin{center}
\begin{tikzpicture}[auto,>=latex,minimum width=.5cm]
  \node[] (input0) {$s_0$};
  \node[below of=input0,node distance=.3cm] (dummy) {};
  \node[below of=dummy,node distance=.3cm] (input1) {$s_1$};
  \node[block, right of=dummy] (T) {$T$};
  \node[right of=T] (output) {$s$};
  \draw[->] (input0) -- (T);
  \draw[->] (input1) -- (T);
  \draw[->] (T) -- (output);
\end{tikzpicture}
\vspace{-.2cm}
\end{center}

We generally think of streams as sequences of \emph{small} values,
and we will use them in this way.
However, we make a leap of imagination and also treat a \emph{whole database} as a stream value.
What is a stream of databases?  It is a \emph{sequence of database
snapshots}.  We model the time-evolution of a database $DB$ as a
stream $DB \in \stream{SCH}$, where $SCH$ is the database schema.
Time is not the wall-clock time, but essentially a counter
of the \emph{sequence of transactions} applied to the database. 
Since transactions are linearizable,
they have a total order, which defines a linear time $t$ dimension:
the value of the stream $DB[t]$ is the snapshot of the 
database contents after $t$ transactions have been applied.  We assume that $DB[0] = 0$, i.e., the database starts empty.

Database transactions also form a stream $T$, a stream of \emph{changes},
or \emph{deltas} that are applied to our database.   
The database snapshot at time $t$ is the cumulative result of applying all 
transactions in the sequence up to $t$: $DB[t] = \sum_{i \leq t} T[i] \defn \I(T)[t]$ (we make the notion of ``addition'' precise later.).
The operation of adding up all changes is \emph{stream integration}.
The following diagram expresses this relationship using the $\I$ operator for 
stream integration:

\begin{center}
\begin{tikzpicture}[auto,>=latex,minimum width=.5cm]
  \node[] (input) {$T$};
  \node[block, right of=input] (I) {$\I$};
  \node[right of=I] (output) {$DB$};
  \draw[->] (input) -- (I);
  \draw[->] (I) -- (output);
\end{tikzpicture}
\end{center}

Conversely, we can say that transactions are the \emph{changes} of a database, and write
$T = \D(DB)$, or $T[t] = DB[t] - DB[t-1]$.  This is the definition of stream differentiation, denoted by $\D$; this 
operation computes the changes of a stream, and is the 
inverse of stream integration.  \secref{sec:streams}
precisely defines streams, integration and differentiation, and analyzes
their properties. 

Let us apply these concepts to view maintenance.
Consider a database $DB$ and a query $Q$ defining a view $V$ as a function
of a database snapshot $V = Q(DB)$.  Corresponding to the stream of database snapshots
$DB$ we have a \emph{stream of view snapshots}: $V[t]$ is the
view's contents after the $t$-th transaction has been applied.  We show this
relationship using the following diagram:

\begin{center}
\begin{tikzpicture}[auto,>=latex,minimum width=.5cm]
  \node[] (input) {$DB$};
  \node[block, right of=input] (I) {$\lift{Q}$};
  \node[right of=I] (output) {$V$};
  \draw[->] (input) -- (I);
  \draw[->] (I) -- (output);
\end{tikzpicture}
\end{center}

The symbol $\lift{Q}$ (the ``lifting'' of $Q$) shows that the query $Q$ is applied
independently to every element of the stream of database snapshots $DB$.  
$\lift{Q}$ is a ``streaming query'' since it operates on a stream of values.
The incremental view maintenance problem requires an algorithm to  
compute the stream $\Delta V$ of \emph{changes} of the view $V$, i.e., $\D(V)$,
as a function of the stream $T$.  
By chaining these definitions together we get the following \defined{fundamental equation}
of the view maintenance problem: $\Delta V = \D(\lift{Q}(DB)) = \D(\lift{Q}(\I(T)))$,
graphically shown as:

\begin{center}
\begin{tikzpicture}[auto,>=latex,minimum width=.5cm]
  \node[] (input) {$T$};
  \node[block, right of=input] (I) {$\I$};
  \node[block, right of=I, node distance=1.5cm] (Q) {$\lift{Q}$};
  \node[block, right of=Q] (D) {$\D$};
  \node[right of=D] (output) {$\Delta V$};
  \draw[->] (input) -- (I);
  \draw[->] (I) -- node (db) {$DB$} (Q);
  \draw[->] (Q) -- (D);
  \draw[->] (D) -- (output);
\end{tikzpicture}
\end{center}

This definition can be generalized to more general streaming queries
$S: \stream{A} \to \stream{B}$ that are richer than lifted pointwise queries $Q$.
The incremental version of streaming query $S$ is denoted by $\inc{S}$ and is defined
according to the above equation, which can also be written as: $\inc{S} = \D \circ S \circ \I$. 

It is generally assumed that the changes to a dataset are much smaller than
the dataset itself; thus, computing on streams of changes may
produce significant performance benefits.

Applying the query \emph{incrementalization} operator $S \mapsto \inc{S}$ constructs
a query that computes directly on changes; however, the resulting query is no more
efficient than a query that computes on the entire dataset, because it uses an
integration operator to reconstitute the full dataset.
\secref{sec:incremental} shows how algebraic properties of the $\inc{\cdot}$ operator 
are used to optimize the implementation of $\inc{S}$:

\begin{enumerate}[nosep, leftmargin=0pt, itemindent=0.5cm, label=\textbf{(\arabic*)}]
\item The first property is that many classes of primitive operations have very efficient incremental
versions.  In particular, linear queries have the property $Q = \inc{Q}$.  Almost all
relational and Datalog queries are based on linear operators.  Thus, the incremental version
of such queries can be computed in time proportional to the size of the changes.  Bilinear 
operators (such as joins) have a more complex implementation, which nevertheless still performs 
work proportional to the size of the changes, but require storing an amount of data proportional
to the size of the relations.

\item The second key property is the chain rule:
$\inc{(S_1 \circ S_2)} = \inc{S_1} \circ \inc{S_2}$.  This rule gives the incremental
version of a complex query as a composition of incremental versions of its components.
It follows that we can implement any incremental query as a composition of primitive
incremental queries, \emph{all of which perform work proportional to the size of the changes}.
\end{enumerate}

Armed with this general theory of incremental computation, in \secref{sec:relational}  
we show how to model relational queries in \dbsp.  This immediately gives
us a general algorithm to compute the incremental version of any relational query.
These results are well-known, but they are cleanly modeled by \dbsp.

Applying DBSP to recursive queries requires extending this computational model.  
In \secref{sec:recursion} we introduce two additional operators: $\delta_0$
creates a stream from a scalar value, and $\int$ creates a scalar value from a stream.  
These operators can be used to implement computations with \code{while} loops.
So, in addition to modelling changing inputs and database, we also
use streams as a model for sequences of \emph{consecutive values of loop
iteration variables}.  With this addition \dbsp becomes rich enough to implement 
recursive queries.  \secref{sec:datalog} shows how stratified recursive
Datalog programs with negation can be implemented in \dbsp.  

In \secref{sec:nested} we use \dbsp to model computations on nested streams, where each 
value of a stream is another stream.  This allows 
us to define \emph{incremental streaming computations for recursive programs}.  As a 
consequence we derive a universal algorithm for incrementalizing arbitrary streaming Datalog programs.

\dbsp is a \textbf{simple} language: 
the basic  \dbsp streaming model is built essentially from two elementary mathematical operators: lifting $\lift$ and delay $\zm$.
The nested streams model adds two additional operators, $\delta_0$ for stream construction and $\int$ for destruction.


\dbsp is also \textbf{expressive}: for example, it is more powerful than stratified Datalog.  
\dbsp can also describe streaming window queries, or queries on nested relations (such as grouping),
and non-monotone recursive queries. 
We discuss briefly the application of \dbsp to richer languages in \secref{sec:extensions}.

\begin{comment}
    \item[abstract:] the two core concepts of 
    (1) streaming computations, and (2) incremental computations
    are completely abstract, since they work over any algebraic group structure.  
    By instantiating the group structure with the real numbers $\mathbb{R}$, 
    we get traditional signal processing systems.  By substituting
    for the group structure \zrs we obtain streaming incremental computations over databases. 
    \item[intuitive:] we often find that the connections of \dbsp to signal processing 
    make it easier to understand.  It is easier to think about the theory as
    computing over real values, samples of a continuous signal.  Moreover, the two
    stream operators central to \dbsp, integration and differentiation,
    have very intuitive meanings in the world of signals.
    \item[visual:] In many circumstances stream circuit diagrams 
    are easier to read and comprehend than the corresponding mathematical formulas.
    \item[economic:] the simple streaming model is built essentially from
    two elementary mathematical operators: lifting $\lift$ and delay $\zm$.  The nested
    streams model adds two additional operators, $\delta_0$ for stream construction and
    $\int$ for destruction.
\end{comment}

This paper omits most proofs; the full proofs are available in an expansive  
companion technical report~\cite{tr}.

This paper makes the following contributions:
\begin{enumerate}[nosep, leftmargin=0pt, itemindent=0.5cm, label=\textbf{(\arabic{*})}]
    \item It defines \dbsp, a small language for streaming computation, which
    nonetheless can express nested non-monotonic recursion;
    \item It provides an algorithm incrementalizing all \dbsp programs;
    \item For fragments of \dbsp corresponding to the relational algebra and stratified-monotonic 
    Datalog, 
    the automatic incrementalization algorithm provides results matching state-of-the-art approaches.
    Moreover, our approach also applies to more powerful languages, such as while-relational and non-monotonic 
    Datalog~\cite{Abiteboul-book95}.
    \item It develops a formal, sound foundation for the manipulation of streaming and incremental computations,
    which allows one to reason formally about program transformations and design new efficient implementations.
    \item \dbsp can express both streaming and incremental computation models in a single framework.  We regard this unification as a significant contribution.
\end{enumerate}